\section{Database Architecture}

ByteVault relies on a distributed PostgreSQL architecture designed to simulate a multi-branch banking environment. This architecture enforces strict separation of data at the storage level while allowing secure, atomic cross-branch transactions.

\subsection{Technology Justification: Why PostgreSQL?}
While NoSQL databases (like MongoDB or Cassandra) excel in horizontal scalability and schema flexibility, ByteVault explicitly mandates \textbf{PostgreSQL} for its ledger foundation. This choice is rooted in the mathematical requirements of financial systems:
\begin{itemize}
    \item \textbf{Strict ACID Compliance}: Financial ledgers cannot tolerate "eventual consistency." If money is deducted from Account A, it must be instantaneously and atomically credited to Account B. PostgreSQL provides rock-solid ACID guarantees natively.
    \item \textbf{Relational Integrity}: The double-entry ledger requires complex foreign-key constraints (e.g., ensuring a \texttt{journal\_line} maps to a valid \texttt{ledger\_account}). Relational databases enforce this structural integrity at the storage layer.
    \item \textbf{Advanced Distributed Features}: PostgreSQL natively supports the Two-Phase Commit (2PC) protocol and Foreign Data Wrappers (FDW), which are essential for coordinating cross-branch transactions in ByteVault.
\end{itemize}

\subsection{Multi-Container Setup}
The system utilizes two distinct PostgreSQL instances, defined in \texttt{docker-compose.yml}:
\begin{itemize}
  \item \textbf{Main DB (Panji Branch)}: \texttt{branch\_a\_db} listening on internal port 5432 (mapped to 55432).
  \item \textbf{Sub DB (Mapusa Branch)}: \texttt{branch\_b\_db} listening on internal port 5432 (mapped to 55433).
\end{itemize}

\textbf{Why it is there:} This multi-database setup models real-world core banking systems where regional branches might operate on physically separate infrastructure to ensure data sovereignty and reduce blast radius during failures.

\subsection{PostgreSQL Foreign Data Wrapper (FDW)}
To facilitate cross-branch operations without the backend needing to manually coordinate dual database connections for simple queries, ByteVault uses \texttt{postgres\_fdw}. 

The Main DB (Panji) mounts a subset of the Sub DB (Mapusa) tables into a specialized schema named \texttt{fdw\_sub}.

\begin{lstlisting}[style=sql]
-- Snippet from docker/init.sql
CREATE SERVER IF NOT EXISTS sub_branch
  FOREIGN DATA WRAPPER postgres_fdw
  OPTIONS (host 'branch_b_db', port '5432', dbname 'branch_b_db');

CREATE USER MAPPING IF NOT EXISTS FOR CURRENT_USER SERVER sub_branch
  OPTIONS (user 'admin_b', password 'secure_password_b');

IMPORT FOREIGN SCHEMA public
  LIMIT TO (branches, users, employees, accounts)
  FROM SERVER sub_branch INTO fdw_sub;
\end{lstlisting}

\textbf{Why it is there:} When a user in the Main branch initiates a transfer to the Sub branch, the backend can effortlessly validate the destination account by querying \texttt{SELECT id FROM fdw\_sub.accounts WHERE account\_number = \$1}. This abstracts the network boundary at the database layer.

\subsection{Two-Phase Commit (2PC)}
Cross-branch financial transfers cannot be handled by standard local transactions. If the Main DB deducts funds, but the Sub DB crashes before crediting them, money is destroyed. 

To solve this, ByteVault uses PostgreSQL's Distributed Two-Phase Commit protocol.

\begin{lstlisting}[style=bash]
# PostgreSQL Configuration (docker/custom.conf)
max_prepared_transactions = 100
\end{lstlisting}

The 2PC workflow ensures Atomicity across network boundaries:
\begin{enumerate}
    \item \textbf{Prepare Phase}: The backend instructs both Main and Sub databases to \texttt{PREPARE TRANSACTION 'tx-uuid'}. Both databases lock the required rows and write the pending transaction to disk, promising they *can* commit.
    \item \textbf{Commit Phase}: If both respond successfully, the backend issues \texttt{COMMIT PREPARED 'tx-uuid'} to both.
    \item \textbf{Rollback Phase}: If either DB fails during the Prepare phase, the backend issues \texttt{ROLLBACK PREPARED} to cancel the operation universally.
\end{enumerate}

\textbf{Why it is there:} 2PC guarantees that distributed systems achieve consensus. A transaction is either 100\% committed on both branches, or 0\% committed, eliminating the risk of partial funds transfer.

\subsection{Schema Design}
The relational schema is highly normalized:
\begin{itemize}
    \item \textbf{Identity}: \texttt{users} (Customers) and \texttt{employees} (Staff).
    \item \textbf{Security}: \texttt{employee\_credentials} (Isolates password hashes from public profile data).
    \item \textbf{Core Banking}: \texttt{accounts} and \texttt{account\_holds} (For temporarily reserving funds during pending transfers).
    \item \textbf{Workflow}: \texttt{transfer\_requests} (Maker/Checker approval pipelines).
    \item \textbf{Double-Entry Ledger}: \texttt{journal\_entries} and \texttt{journal\_lines}.
\end{itemize}
